\documentclass[11pt, a4paper]{article}
\usepackage[utf8]{inputenc}
\usepackage{geometry}
\usepackage{booktabs}
\usepackage{graphicx}
\usepackage{hyperref}
\usepackage{amsmath}
\usepackage{enumitem}
\usepackage{titlesec}
\usepackage{xcolor}

\geometry{a4paper, margin=1in}
\setlength{\parindent}{0em}
\setlength{\parskip}{1em}

\begin{document}

\begin{titlepage}
    \centering
    \vspace*{3cm}
    
    {\Huge \textbf{Milestone 2: Data Engineering \& Validation}}\\[1.5cm]
    {\huge \textbf{ReSpace-AI Project}}\\[2.5cm]
    
    {\Large \textbf{Group Members:}}\\[0.5cm]
    {\large
    \renewcommand{\arraystretch}{1.5}
    \begin{tabular}{lc}
        \textbf{Fatima Mansoor} & CS22015 \\
        \textbf{Muneeb Ahmed} & CS22048 \\
        \textbf{Zobia Naseem} & CS22031 \\
        \textbf{Maham Ahmed} & CS22032 \\
    \end{tabular}
    }\\[3cm]
    
    \vfill
    
    {\Large \today}
\end{titlepage}

\tableofcontents
\newpage

\section{Introduction}
This report outlines the end-to-end data pipeline implemented for our AI-based interior design generator. It covers the transition from raw images to a standardized, machine-learning-ready dataset, encompassing annotation, preprocessing, augmentation, exploratory analysis, and dataset splitting.

\section{Annotation Process}
The annotation methodology prioritizes structured, multi-dimensional metadata extraction. We built a custom GUI application (`customtkinter`-based) to streamline labeling and maintain strict standard schemas.

\begin{figure}[h!]
    \centering
    % \includegraphics[width=0.8\textwidth]{placeholder_gui.png} % Replace 'placeholder_gui.png' with your GUI screenshot
    \framebox(300,150){\Large GUI Application Screenshot}
    \caption{CustomTkinter GUI used for Data Annotation}
    \label{fig:gui_annotation}
\end{figure}

\begin{itemize}
    \item \textbf{Metadata Capture:} Images are systematically tagged across several domain-specific dimensions: \textit{Room Type}, \textit{Color Theme}, \textit{Use Case}, \textit{Lighting}, \textit{Color Palette}, and a multi-select array of \textit{Furniture}.
    \item \textbf{Collaborative Workflow:} The annotations and references are synchronized with a Supabase PostgreSQL backend. To prevent duplicated effort among reviewers, a status-tracking system (`submitted`, `discarded`) restricts an image to be processed solely once.
    \item \textbf{Methodology Justification:} Structuring the qualitative nuances into defined tabular attributes transforms the subjective image context into structured text prompts, maximizing fine-tuning efficacy and categorical balance. 
\end{itemize}

\section{Preprocessing Pipeline}
To ensure all machine-ready data is uniform and of high qualitative fidelity, a rigid filtering and standardization pipeline was implemented prior to split and augmentation:

\begin{itemize}
    \item \textbf{Resizing:} Images were structurally scaled to a high-resolution square of $1024 \times 1024$ pixels. 
    \textit{Justification:} A standardized square resolution is an architectural requirement for our downstream generative networks (e.g., Stable Diffusion/Transformers).
    \item \textbf{Normalization:} Pixel intensities were normalized uniformly onto the interval $[0, 1]$ directly mapping into an `float32` datatype representation. 
    \textit{Justification:} Min-max scaling accelerates and stabilizes gradient-based optimization during training periods.
    \item \textbf{Dynamic Quality Reduction (Noise/Blur/Spots):}
        \begin{itemize}
            \item \textbf{Blur filter:} OpenCV's Laplacian variance filter is applied. Images returning a variance threshold below 250.0 were automatically discarded as blurry.
            \item \textbf{Exposure checking:} Used binary thresholding to trace bright-spot anomalies. Unnatural exposure spots were filtered out if binary variance fell beyond a natural acceptable band ($5000 < \text{variance} < 8500$). 
        \end{itemize}
\end{itemize}

\section{Exploratory Data Analysis (EDA)}
Using Seaborn and Pandas, comprehensive statistical analysis was carried out to audit category representations.

\begin{figure}[h!]
    \centering
    % \includegraphics[width=0.9\textwidth]{placeholder_eda_histograms.png} % Replace 'placeholder_eda_histograms.png' with your histogram
    \framebox(350,150){\Large Class Distribution Histograms}
    \caption{Distribution of Room Types, Color Themes, and Lighting}
    \label{fig:eda_histograms}
\end{figure}

\begin{itemize}
    \item \textbf{Class Distribution Histograms:} Univariate distributions revealed our underlying sampling biases (e.g., specific \textit{Lighting} configurations vs. \textit{Color Palettes}). Entropy balances natively visualized through Shannon Entropy ratios helped trace overrepresented subsets, guiding strategic data resampling.
    \item \textbf{Cross-Category Dynamics:} Multi-variate relationship associations (Cramér's V correlation measurements) proved independence between structural metrics like ``Furniture Item Counts'' and varying ``Room Types.'' 
    \item \textbf{Image Resolution Statistics:} As a byproduct of our deterministic preprocessing pipeline, 100\% of the underlying analytical dataset structurally aligns to an exact resolution configuration ($1024 \times 1024 \text{ px}$), thus generating a uniform bounding box static profile that renders resolution-based standard-deviation negligible prior to entry into spatial augmentations.
\end{itemize}

\section{Data Augmentation Strategy}
To artificially scale dataset volume while building robustness against typical spatial and photography constraints, we generated two definitive programmatic variants per image:

\begin{itemize}
    \item \textbf{Random Scaling \& Cropping:} The original image is randomly resized using a relative transformation factor (between 0.7x and 1.0x), then randomly cropped back to the base architectural requirement ($1024 \times 1024$).
    \item \textbf{Horizontal Flipping:} Geometrically flipping the perspective layouts around the Y-axis.
\end{itemize}

\textbf{Real-World Justification:} 
These strategies strictly address real-world variations typical inside real estate cataloging. \textbf{Random crops} simulate different camera focal lengths, field-of-views, and interior framings (i.e. focusing solely on a desk arrangement versus an overarching wide-angle capture of a room). \textbf{Horizontal flips} imitate varied floor-plan layouts and mirrored natural lighting entries. These augmentations mathematically expand representation space without altering or decaying the precise semantic attributes described in text prompts. 

\section{Data Split Documentation}
To accurately frame real generalizability, dataset division is rigorously isolated via a randomized subset allocation.

\begin{table}[h!]
    \centering
    \renewcommand{\arraystretch}{1.3}
    \begin{tabular}{@{}lcc@{}}
        \toprule
        \textbf{Dataset Partition} & \textbf{Allocation (\%)} & \textbf{Purpose} \\
        \midrule
        Training Set & 80\% & Model training and anchor for geometric augmentations. \\
        Test Set & 20\% & Holdout portion for unbiased generative evaluation. \\
        \bottomrule
    \end{tabular}
    \caption{Dataset Partitioning Strategy}
    \label{tab:data_split}
\end{table}

\begin{itemize}
    \item \textbf{Training Set (80\%):} The primary partition strictly isolated. Augmented generations are derived \textit{only} from subjects existing inside this training boundary to preclude geometric data-leakage.
    \item \textbf{Test Set (20\%):} Kept wholly untouched and purely preserved from both upscaling augmentations and training observation logic.
    \item \textbf{Validation Set:} Typically allocated as an inner-fold (for instance, a 90/10 inner split inside the 80\% Training partition dynamically allocated during modeling iterations) facilitating epoch-based hyperparameter tuning without exposing test cases.
\end{itemize}

\end{document}
